\documentclass[12pt,a4paper]{article}

%\usepackage[margin=1in]{geometry}
\usepackage[a4paper,margin=1in]{geometry}
\usepackage{graphicx}
\usepackage{setspace}
\usepackage{titlesec}
\usepackage{xcolor}
\usepackage{listings}

\lstset{
  language=SQL,
  basicstyle=\ttfamily\small,
  frame=single,
  breaklines=true,
  showstringspaces=false,
  columns=fullflexible
}

\begin{document}

\section*{Chapter 1}

\subsection*{Question 1}
The following SELECT statement executes successfully:
\begin{lstlisting}
    SELECT  last_name,
            job_id,
            salary AS Sal
    FROM    employees;
\end{lstlisting}

\textbf{Answer:} 

Yes, the statement executes successfully.

\subsection*{Question 2}
The following SELECT statement executes successfully:
\begin{lstlisting}
    SELECT  *
    FROM    job_grades;
\end{lstlisting}

\textbf{Answer:}

No, the statement does not execute successfully because the table \texttt{job\_grades} does not exist (in the given schema/environment).

\subsection*{Question 3}
There are four coding errors in the following statement. Can you identify them?
\begin{lstlisting}
    SELECT  employee_id, last_name
        sal x 12 ANNUAL SALARY
    FROM    employees;
\end{lstlisting}

\textbf{Answer:}
    \textbf{Errors:}
        \begin{enumerate}
            \item Missing comma after \texttt{last\_name}.
            \item \texttt{sal} is not a valid column name, it should be \texttt{salary}.
            \item \texttt{x} is invalid for multiplication, we need to use \texttt{*}.
            \item Alias \texttt{ANNUAL SALARY} must be enclosed in double quotes: \texttt{AS "ANNUAL SALARY"}.
        \end{enumerate}

\textbf{Correct query:}
\begin{lstlisting}
    SELECT  employee_id, last_name,
    salary  * 12 AS "ANNUAL SALARY"
    FROM    employees;
\end{lstlisting}

\subsection*{Question 4}
Your first task is to determine the structure of the \texttt{DEPARTMENTS} table and its contents.

\textbf{Answer:}
    \begin{lstlisting}
        DESCRIBE    departments;
        SELECT  *
        FROM    departments;
    \end{lstlisting}

\subsection*{Question 5}
The HR department wants a query to display the last name, job ID, hire date,and employee ID 
for each employee, with the employee ID appearing first. 
Provide an alias STARTDATE for the \texttt{HIRE\_DATE} column. 
Save your SQL statement to a file named \texttt{lab\_01\_05.sql} so that you can 
dispatch this file to the HR department.

\textbf{Answer:}
    \begin{lstlisting}
        SELECT  employee_id, last_name, job_id, 
        hire_date AS STARTDATE
        FROM employees;
    \end{lstlisting}

\subsection*{Question 6}
Test your query in the \texttt{lab\_01\_05.sql} file to ensure that it runs correctly.
Note: After you have executed the query, make sure that you do not enter 
your next query in the same worksheet. Open a new worksheet.

\textbf{Answer:}
    \begin{lstlisting}
        To open a new worksheet: CTRL + SHIFT + N
    \end{lstlisting}

\subsection*{Question 7}
The HR department wants a query to display all unique job IDs from the EMPLOYEES table.

\textbf{Answer:}
    \begin{lstlisting}
        SELECT DISTINCT job_id
        FROM employees;
    \end{lstlisting}

\subsection*{Question 8}
The HR department wants more descriptive column headings for its report on employees. 
Copy the statement from \texttt{lab\_01\_05.sql} to a new SQL Worksheet. Name the column headings 
\texttt{Emp \#}, \texttt{Employee}, \texttt{Job}, and \texttt{Hire Date}, respectively. Then run your query again.

\textbf{Answer:}
    \begin{lstlisting}
        SELECT employee_id AS "Emp #",
        last_name AS "Employee",
        job_id AS "Job",
        hire_date AS "Hire Date"
        FROM employees;
    \end{lstlisting}

\newpage
\subsection*{Question 9}
The HR department has requested a report of all employees and their job IDs. Display the last 
name concatenated with the job ID (separated by a comma and space) and name the column 
\texttt{Employee and Title}.

\textbf{Answer:}
    \begin{lstlisting}
        SELECT last_name || ', ' || job_id AS "Employee and Title"
        FROM employees;
    \end{lstlisting}

\subsection*{Question 10}
To familiarize yourself with the data in the \texttt{EMPLOYEES} table, create a query to display all the 
data from that table. Separate each column output by a comma. Name the column title 
\texttt{THE\_OUTPUT}.

\textbf{Answer:}
    \begin{lstlisting}
        SELECT employee_id || ',' || first_name || ',' || last_name || ',' ||
        email || ',' || phone_number || ',' || job_id || ',' ||
        manager_id || ',' || hire_date || ',' || salary || ',' ||
        commission_pct || ',' || department_id AS "THE_OUTPUT"
        FROM employees;
    \end{lstlisting}

\section*{Chapter 2}

\subsection*{Question 1}
Because of budget issues, the HR department needs a report that displays the last name and 
salary of employees who earn more than \$12,000. Save your SQL statement as a file named 
\texttt{lab\_02\_01.sql}. Run your query.

\textbf{Answer:}
    \begin{lstlisting}
        SELECT last_name, salary
        FROM employees
        WHERE salary > 12000;
    \end{lstlisting}

\subsection*{Question 2}
Open a new SQL Worksheet. Create a report that displays the last name and department number 
for employee number 176. Run the query.

\textbf{Answer:}
    \begin{lstlisting}
        SELECT last_name, department_id
        FROM EMPLOYEES
        WHERE employee_id = 176;
    \end{lstlisting}

\subsection*{Question 3}
The HR department needs to find high-salary and low-salary employees. 
Modify \texttt{lab\_02\_01.sql} to display the last name and salary for any employee whose salary is not in 
the range of \$5,000 to \$12,000. Save your SQL statement as \texttt{lab\_02\_03.sql}.

\textbf{Answer:}
    \begin{lstlisting}
        SELECT last_name, salary
        FROM EMPLOYEES
        WHERE salary NOT BETWEEN 5000 AND 12000;
    \end{lstlisting}

\subsection*{Question 4}
Create a report to display the last name, job ID, and hire date for employees with the last names 
of Matos and Taylor. Order the query in ascending order by the hire date.

\textbf{Answer:}
    \begin{lstlisting}
        SELECT last_name, job_id, hire_date
        FROM employees
        WHERE last_name IN ('Matos', 'Taylor')
        ORDER BY hire_date ASC;
    \end{lstlisting}

\subsection*{Question 5}
Display the last name and department ID of all employees in departments 20 or 50 in ascending 
alphabetical order by name.

\textbf{Answer:}
    \begin{lstlisting}
        SELECT last_name, department_id
        FROM employees
        WHERE department_id IN (20, 50)
        ORDER BY last_name ASC;
    \end{lstlisting}

\subsection*{Question 6}
Modify \texttt{lab\_02\_03.sql} to display the last name and salary of employees who earn between 
\$5,000 and \$12,000, and are in department 20 or 50. Label the columns Employee and 
Monthly Salary, respectively. Resave \texttt{lab\_02\_03.sql} as \texttt{lab\_02\_06.sql}. Run the 
statement in \texttt{lab\_02\_06.sql}.

\textbf{Answer:}
    \begin{lstlisting}
        SELECT last_name AS "Employee",
        salary AS "Monthly Salary"
        FROM employees
        WHERE salary BETWEEN 5000 AND 12000
        AND department_id IN (20, 50);
    \end{lstlisting}

\subsection*{Question 7}
The HR department needs a report that displays the last name and hire date for all employees 
who were hired in 1994.

\textbf{Answer:}
    \begin{lstlisting}
        SELECT last_name, hire_date
        FROM employees
        WHERE TO_CHAR(hire_date, 'YYYY') = '1994';
    \end{lstlisting}

\subsection*{Question 8}
Create a report to display the last name and job title of all employees who do not have a 
manager.

\textbf{Answer:}
    \begin{lstlisting}
        SELECT last_name, job_id
        FROM employees
        WHERE manager_id IS NULL;
    \end{lstlisting}

\subsection*{Question 9}
Create a report to display the last name, salary, and commission of all employees who earn 
commissions. Sort data in descending order of salary and commissions.
Use the column’s numeric position in the ORDER BY clause.

\textbf{Answer:}
    \begin{lstlisting}
        SELECT last_name, salary, commission_pct
        FROM employees
        WHERE commission_pct IS NOT NULL
        ORDER BY 2 DESC, 3 DESC;
    \end{lstlisting}

\subsection*{Question 10}
Members of the HR department want to have more flexibility with the queries that you are 
writing. They would like a report that displays the last name and salary of employees who earn 
more than an amount that the user specifies after a prompt. Save this query to a file named 
\texttt{lab\_02\_10.sql}. If you enter 12000 when prompted, the report displays the following 
results:

\textbf{Answer:}
    \begin{lstlisting}
        SELECT last_name, salary
        FROM employees
        WHERE salary > &amount;
    \end{lstlisting}

\subsection*{Question 11}
The HR department wants to run reports based on a manager. Create a query that prompts the 
user for a manager ID and generates the employee ID, last name, salary, and department for 
that manager’s employees. The HR department wants the ability to sort the report on a selected 
column. You can test the data with the following values:

\textbf{Answer:}
    \begin{lstlisting}
        SELECT employee_id, last_name, salary, department_id
        FROM employees
        WHERE manager_id = &manager_id
        ORDER BY &sort_column;
    \end{lstlisting}

\subsection*{Question 12}
Display all employee last names in which the third letter of the name is “a.”

\textbf{Answer}
    \begin{lstlisting}
        SELECT last_name
        FROM employees
        WHERE last_name LIKE '__a%';
    \end{lstlisting}

\subsection*{Question 13}
Display the last names of all employees who have both an “a” and an “e” in their last name.

\textbf{Answer:}
    \begin{lstlisting}
        SELECT last_name
        FROM employees
        WHERE last_name LIKE '%a%'
        AND last_name LIKE '%e%';
    \end{lstlisting}

\subsection*{Question 14}
Display the last name, job, and salary for all employees whose jobs are either those of a sales 
representative or of a stock clerk, and whose salaries are not equal to \$2,500, \$3,500, or \$7,000.

\textbf{Answer:}
    \begin{lstlisting}
        SELECT last_name, job_id, salary
        FROM employees
        WHERE job_id IN ('SA_REP', 'ST_CLERK')
        AND salary NOT IN (2500, 3500, 7000);
    \end{lstlisting}

\subsection*{Question 15}
Modify \texttt{lab\_02\_06.sql} to display the last name, salary, and commission for all employees 
whose commission is 20\%. Resave \texttt{lab\_02\_06.sql} as \texttt{lab\_02\_15.sql}. Rerun the 
statement in \texttt{lab\_02\_15.sql}.

\textbf{Answer:}
    \begin{lstlisting}
        SELECT last_name AS "Employee",
        salary AS "Monthly Salary",
        commission_pct
        FROM employees
        WHERE commission_pct = 0.2;
    \end{lstlisting}

\section*{Chapter 3}

\subsection*{Question 1}
Write a query to display the system date. Label the column as Date.
Note: If your database is remotely located in a different time zone, the output will be the date 
for the operating system on which the database resides.

\textbf{Answer:}
    \begin{lstlisting}
        SELECT SYSDATE AS "Date"
        FROM dual;
    \end{lstlisting}

\subsection*{Question 2}
The HR department needs a report to display the employee number, last name, salary, and 
salary increased by 15.5\% (expressed as a whole number) for each employee. Label the column 
\texttt{New Salary}. Save your SQL statement in a file named \texttt{lab\_03\_02.sql}.

\textbf{Answer:}
    \begin{lstlisting}
        SELECT employee_id,
        last_name,
        salary,
        ROUND(salary * 1.155) AS "New Salary"
        FROM employees;
    \end{lstlisting}

\subsection*{Question 3}
Run your query in the \texttt{lab\_03\_02.sql} file.
\textbf{Answer:}
        We can run the following command\: \texttt{@lab\_03\_02.sql}

\subsection*{Question 4}
Modify your query \texttt{lab\_03\_02.sql} to add a column that subtracts the old salary from the 
new salary. Label the column \texttt{Increase}. Save the contents of the file as \texttt{lab\_03\_04.sql}. 
Run the revised query.

\textbf{Answer:}
    \begin{lstlisting}
        SELECT employee_id,
        last_name,
        salary,
        ROUND(salary * 1.155) AS "New Salary",
        ROUND((salary * 1.155) - salary) AS "Increase"
        FROM employees;
    \end{lstlisting}

\subsection*{Question 5}
Write a query that displays the last name (with the first letter in uppercase and all the other 
letters in lowercase) and the length of the last name for all employees whose name starts with 
the letters “J,” “A,” or “M.” Give each column an appropriate label. Sort the results by the 
employees’ last names.

\textbf{Answer:}
    \begin{lstlisting}
        SELECT INITCAP(last_name) AS "Name",
        LENGTH(last_name) AS "Length"
        FROM employees
        WHERE last_name LIKE 'J%'
        OR last_name LIKE 'A%'
        OR last_name LIKE 'M%'
        ORDER BY last_name;
    \end{lstlisting}

Rewrite the query so that the user is prompted to enter a letter that the last name starts with. For 
example, if the user enters “H” (capitalized) when prompted for a letter, then the output should 
show all employees whose last name starts with the letter “H.”

\textbf{Answer:}
    \begin{lstlisting}
        SELECT INITCAP(last_name) AS "Name",
        LENGTH(last_name) AS "Length"
        FROM employees
        WHERE last_name LIKE '&start_letter%'
        ORDER BY last_name;
    \end{lstlisting}

Modify the query such that the case of the entered letter does not affect the output. The entered 
letter must be capitalized before being processed by the \texttt{SELECT} query.

\textbf{Answer:}
    \begin{lstlisting}
        SELECT INITCAP(last_name) AS "Name",
        LENGTH(last_name) AS "Length"
        FROM employees
        WHERE last_name LIKE UPPER('&start_letter') || '%'
        ORDER BY last_name;
    \end{lstlisting}

\subsection*{Question 6}
The HR department wants to find the duration of employment for each employee. For each 
employee, display the last name and calculate the number of months between today and the 
date on which the employee was hired. Label the column as \texttt{MONTHS\_WORKED}. Order your 
results by the number of months employed. Round the number of months up to the closest 
whole number.
Note: Because this query depends on the date when it was executed, the values in the 
\texttt{MONTHS\_WORKED} column will differ for you.

\textbf{Answer:}
    \begin{lstlisting}
        SELECT last_name,
        CEIL(MONTHS_BETWEEN(SYSDATE, hire_date)) AS "MONTHS_WORKED"
        FROM employees
        ORDER BY "MONTHS_WORKED";
    \end{lstlisting}

\subsection*{Question 7}
Create a query to display the last name and salary for all employees. Format the salary to be 15 
characters long, left-padded with the \$ symbol. Label the column as \texttt{SALARY}.

\textbf{Answer:}
    \begin{lstlisting}
        SELECT last_name,
        LPAD(salary, 15, '$') AS "SALARY"
        FROM employees;
    \end{lstlisting}

\subsection*{Question 8}
Create a query that displays the first eight characters of the employees’ last names and indicates 
the amounts of their salaries with asterisks. Each asterisk signifies a thousand dollars. Sort the 
data in descending order of salary. Label the column as 
\texttt{EMPLOYEES\_AND\_THEIR\_SALARIES}.

\textbf{Answer:}
    \begin{lstlisting}
        SELECT SUBSTR(last_name, 1, 8) || ' ' ||
        RPAD('*', FLOOR(salary/1000), '*') AS "EMPLOYEES_AND_THEIR_SALARIES"
        FROM employees
        ORDER BY salary DESC;
    \end{lstlisting}

\subsection*{Question 9}
Create a query to display the last name and the number of weeks employed for all employees in 
department 90. Label the number of weeks column as \texttt{TENURE}. Truncate the number of weeks 
value to 0 decimal places. Show the records in descending order of the employee’s tenure.
Note: The \texttt{TENURE} value will differ as it depends on the date on which you run the query.

\textbf{Answer:}
    \begin{lstlisting}
        SELECT last_name,
        TRUNC(MONTHS_BETWEEN(SYSDATE, hire_date) * 4) AS "TENURE"
        FROM employees
        WHERE department_id = 90
        ORDER BY "TENURE" DESC;
    \end{lstlisting}

\section*{Chapter 4}

\subsection*{Question 1}
Create a report that produces the following for each employee:
\texttt{<employee last name>} earns \texttt{<salary>} monthly 
but wants \texttt{<3 times salary.>}. Label the column \texttt{Dream Salaries}.

\textbf{Answer:}
    \begin{lstlisting}
        SELECT last_name || ' earns $' || TO_CHAR(salary, '99,999.00') ||
        ' monthly but wants $' || TO_CHAR(salary * 3, '99,999.00') AS "Dream Salaries"
        FROM employees;
    \end{lstlisting}

\subsection*{Question 2}
Display each employee’s last name, hire date, and salary review date, which is the first Monday 
after six months of service. Label the column \texttt{REVIEW}. Format the dates to appear in the format 
similar to “Monday, the Thirty-First of July, 2000.”

\textbf{Answer:}
    \begin{lstlisting}
        SELECT last_name,
        hire_date,
        TO_CHAR(NEXT_DAY(ADD_MONTHS(hire_date, 6), 'MONDAY'),
        'fmDay, "the" DDSPTH "of" fmMonth, YYYY') AS "REVIEW"
        FROM employees;
    \end{lstlisting}
    
\subsection*{Question 3}
Display the last name, hire date, and day of the week on which the employee started. Label the 
column DAY. Order the results by the day of the week, starting with Monday.

\textbf{Answer:}
    \begin{lstlisting}
        SELECT last_name,
        hire_date,
        TO_CHAR(hire_date, 'DAY') AS "DAY"
        FROM employees
        ORDER BY TO_CHAR(hire_date, 'D');
    \end{lstlisting}
    
\subsection*{Question 4}
Create a query that displays the employees’ last names and commission amounts. If an employee 
does not earn commission, show “No Commission.” Label the column \texttt{COMM}.

\textbf{Answer:}
    \begin{lstlisting}
        SELECT  last_name,
        NVL(TO_CHAR(salary * commission_pct), 'No Commission') AS "COMM"
        FROM    employees;
    \end{lstlisting}
    
\subsection*{Question 5}
Using the \texttt{DECODE} function, write a query that displays the grade of all employees based on the 
value of the column \texttt{JOB\_ID}, using the following data:
    Job                 Grade
    AD\_PRES             A
    ST\_MAN              B
    IT\_PROG             C
    SA\_REP              D
    ST\_CLERK            E
    None of the above    0

\textbf{Answer:}
    \begin{lstlisting}
        SELECT job_id,
        DECODE(job_id,
              'AD_PRES', 'A',
              'ST_MAN',  'B',
              'IT_PROG', 'C',
              'SA_REP',  'D',
              'ST_CLERK','E',
              '0') AS "GRADE"
        FROM employees;
    \end{lstlisting}
    
\subsection*{Question 6}
Rewrite the statement in the preceding exercise using the CASE syntax.

\textbf{Answer:}
    \begin{lstlisting}
        SELECT job_id,
        CASE job_id
            WHEN 'AD_PRES'  THEN 'A'
            WHEN 'ST_MAN'   THEN 'B'
            WHEN 'IT_PROG'  THEN 'C'
            WHEN 'SA_REP'   THEN 'D'
            WHEN 'ST_CLERK' THEN 'E'
            ELSE '0'
        END AS "GRADE"
        FROM employees;
    \end{lstlisting}
    
\section*{Chapter 5}

\subsection*{Question 1}
Group functions work across many rows to produce one result per group.
True/False

\textbf{Answer:}
    True
    Group functions like \texttt{SUM}, \texttt{AVG}, \texttt{MAX}, \texttt{MIN}, and \texttt{COUNT} operate on sets of rows and return a single result per group.

\subsection*{Question 2}
Group functions include nulls in calculations.
True/False

\textbf{Answer:}
    False
    Most group functions ignore nulls. For example, \texttt{AVG}, \texttt{SUM}, and \texttt{COUNT} skip null values unless you use \texttt{COUNT(*)}.

\subsection*{Question 3}
The \texttt{WHERE} clause restricts rows before inclusion in a group calculation.
True/False

\textbf{Answer:}
    True
    The \texttt{WHERE} clause filters rows before any grouping or aggregation happens.

\subsection*{Question 4}
Find the highest, lowest, sum, and average salary of all employees. Label the columns
as \texttt{Maximum}, \texttt{Minimum}, \texttt{Sum}, and \texttt{Average}, respectively. Round your results to the nearest 
whole number. Save your SQL statement as \texttt{lab\_05\_04.sql}. Run the query.

\textbf{Answer:}
    \begin{lstlisting}
        SELECT ROUND(MAX(salary)) AS "Maximum",
        ROUND(MIN(salary)) AS "Minimum",
        ROUND(SUM(salary)) AS "Sum",
        ROUND(AVG(salary)) AS "Average"
        FROM employees;
    \end{lstlisting}

\subsection*{Question 5}
Modify the query in \texttt{lab\_05\_04.sql} to display the minimum, maximum, sum, and average 
salary for each job type. Resave \texttt{lab\_05\_04.sql} as \texttt{lab\_05\_05.sql}. Run the statement 
in \texttt{lab\_05\_05.sql}.

\textbf{Answer:}
    \begin{lstlisting}
        SELECT job_id,
        ROUND(MAX(salary)) AS "Maximum",
        ROUND(MIN(salary)) AS "Minimum",
        ROUND(SUM(salary)) AS "Sum",
        ROUND(AVG(salary)) AS "Average"
        FROM employees
        GROUP BY job_id;
    \end{lstlisting}
    
\subsection*{Question 6}
Write a query to display the number of people with the same job.

\textbf{Answer:}
    \begin{lstlisting}
        SELECT job_id,
        COUNT(*) AS "COUNT"
        FROM employees
        GROUP BY job_id;
    \end{lstlisting}

Generalize the query so that the user in the HR department is prompted for a job title. Save the script 
to a file named \texttt{lab\_05\_06.sql}. Run the query. Enter \texttt{IT\_PROG} when prompted.

    \begin{lstlisting}
        SELECT job_id,
        COUNT(*) AS "COUNT"
        FROM employees
        WHERE job_id = '&job_title'
        GROUP BY job_id;
    \end{lstlisting}
    
\subsection*{Question 7}
Determine the number of managers without listing them. Label the column as \texttt{Number of Managers}. Hint: Use the \texttt{MANAGER\_ID} column to determine the number of managers.

\textbf{Answer:}
    \begin{lstlisting}
        SELECT COUNT(DISTINCT manager_id) AS "Number of Managers"
        FROM employees;
    \end{lstlisting}
    
\subsection*{Question 8}
Find the difference between the highest and lowest salaries. Label the column \texttt{DIFFERENCE}.

\textbf{Answer:}
    \begin{lstlisting}
        SELECT (MAX(salary) - MIN(salary)) AS "DIFFERENCE"
        FROM employees;
    \end{lstlisting}
    
\subsection*{Question 9}
Create a report to display the manager number and the salary of the lowest-paid employee for 
that manager. Exclude anyone whose manager is not known. Exclude any groups where the 
minimum salary is \$6,000 or less. Sort the output in descending order of salary.

\textbf{Answer:}
    \begin{lstlisting}
        SELECT manager_id,
        MIN(salary) AS "MIN_SALARY"
        FROM employees
        WHERE manager_id IS NOT NULL
        GROUP BY manager_id
        HAVING MIN(salary) > 6000
        ORDER BY MIN(salary) DESC;
    \end{lstlisting}
    
\subsection*{Question 10}
Create a query to display the total number of employees and, of that total, the number of 
employees hired in 1995, 1996, 1997, and 1998. Create appropriate column headings.

\textbf{Answer:}
    \begin{lstlisting}
        SELECT COUNT(*) AS "TOTAL",
        SUM(CASE WHEN TO_CHAR(hire_date, 'YYYY') = '1995' THEN 1 ELSE 0 END) AS "1995",
        SUM(CASE WHEN TO_CHAR(hire_date, 'YYYY') = '1996' THEN 1 ELSE 0 END) AS "1996",
        SUM(CASE WHEN TO_CHAR(hire_date, 'YYYY') = '1997' THEN 1 ELSE 0 END) AS "1997",
        SUM(CASE WHEN TO_CHAR(hire_date, 'YYYY') = '1998' THEN 1 ELSE 0 END) AS "1998"
        FROM employees;
    \end{lstlisting}
    
\subsection*{Question 11}
Create a matrix query to display the job, the salary for that job based on department number, 
and the total salary for that job, for departments 20, 50, 80, and 90, giving each column an 
appropriate heading.

\textbf{Answer:}
    \begin{lstlisting}
        SELECT job_id,
        SUM(CASE WHEN department_id = 20 THEN salary END) AS "Dept 20",
        SUM(CASE WHEN department_id = 50 THEN salary END) AS "Dept 50",
        SUM(CASE WHEN department_id = 80 THEN salary END) AS "Dept 80",
        SUM(CASE WHEN department_id = 90 THEN salary END) AS "Dept 90",
        SUM(salary) AS "Total"
        FROM employees
        WHERE department_id IN (20, 50, 80, 90)
        GROUP BY job_id;
    \end{lstlisting}
    
\end{document}